\sectionkicker{Source-backed technical notes}
\subsection*{生成メディアが研究デモから配布可能な生態系へ――Flamingo、DALL-E 2、Stable Diffusion: Technical Notes}
本欄は記事本文で圧縮した一次資料上の情報を、比較・再検証しやすい形で整理したものである。時系列、確認済みの主張、留意点、一次資料URLを掲載する。
\smallskip
\noindent{\small\textit{Technical Notesでは、一次資料から確認した公開・提供・研究上の事実を記載する。数値・能力評価や提供元・著者による比較表現は、独立再現済みの結果ではなく帰属claimとして扱う。}}\par
\medskip
\subsection*{Theme at a glance}
\addcontentsline{toc}{subsection}{Theme at a glance}
\begin{center}
\footnotesize
\begin{tabularx}{\linewidth}{@{}p{0.20\linewidth}p{0.10\linewidth}p{0.14\linewidth}X@{}}
\toprule
資料 & 位置づけ & 種別 & 時系列 \\
\midrule
Flamingo & 主要資料 & モデル & — \\
DALL-E 2 & 主要資料 & モデル & — \\
Stable Diffusion & 主要資料 & オープンウェイト & — \\
Imagen & 補足資料 & 論文 & — \\
LAION-5B & 補足資料 & 論文 & — \\
DreamBooth & 補足資料 & 論文 & — \\
AudioLM & 補足資料 & 論文 & — \\
Make-A-Video & 補足資料 & 論文 & — \\
DreamFusion & 補足資料 & 論文 & — \\
\bottomrule
\end{tabularx}
\end{center}
\normalsize

\begin{technicalnote}{Flamingo}{主要資料}
\begingroup
% reader-facing Technical Notes break policy
\widowpenalty=10000
\clubpenalty=10000
\displaywidowpenalty=10000
\begin{tabularx}{\linewidth}{@{}>{\bfseries}p{0.22\linewidth}X@{}}
組織 & Google DeepMind \\
種別 & モデル \\
時系列 & — \\
\end{tabularx}
\smallskip
{\bfseries 一次資料から整理したtechnical points}
\begin{itemize}[leftmargin=1.5em,itemsep=0.35em]
\item \textbf{一次情報で確認できる事実}: Google DeepMindの選定済み一次資料では、Flamingoは画像・動画とtextをinterleaveした入力を扱い、少数のtask-specific examplesをpromptとして与えるfew-shot visual-language modelとして設計された。 architectureはpretrained vision modelとlanguage modelを接続し、単一modelで複数のvisual-language taskへ適応することを狙う。
\end{itemize}
\begin{minipage}{\linewidth}
% reader-facing Technical Notes generic boundary/source tail group
\begin{samepage}
{\bfseries 一次資料}
\begingroup\sloppy
\Urlmuskip=0mu plus 2mu\relax
\begin{itemize}[leftmargin=1.5em,itemsep=0.25em]
\item {\scriptsize\url{http://arxiv.org/abs/2204.14198v2}}
\item {\scriptsize\url{https://deepmind.google/blog/tackling-multiple-tasks-with-a-single-visual-language-model/}}
\end{itemize}
\endgroup
\end{samepage}
\end{minipage}
\endgroup
\end{technicalnote}
\medskip

\begin{technicalnote}{DALL-E 2}{主要資料}
\begingroup
% reader-facing Technical Notes break policy
\widowpenalty=10000
\clubpenalty=10000
\displaywidowpenalty=10000
\begin{tabularx}{\linewidth}{@{}>{\bfseries}p{0.22\linewidth}X@{}}
組織 & OpenAI \\
種別 & モデル \\
時系列 & — \\
\end{tabularx}
\smallskip
{\bfseries 一次資料から整理したtechnical points}
\begin{itemize}[leftmargin=1.5em,itemsep=0.35em]
\item \textbf{一次情報で確認できる事実}: OpenAIの選定済み一次資料では、対象event近傍の一次資料から autoregressive modeling を確認できる。
\end{itemize}
\begin{minipage}{\linewidth}
% reader-facing Technical Notes generic boundary/source tail group
\begin{samepage}
{\bfseries 一次資料}
\begingroup\sloppy
\Urlmuskip=0mu plus 2mu\relax
\begin{itemize}[leftmargin=1.5em,itemsep=0.25em]
\item {\scriptsize\url{http://arxiv.org/abs/2204.06125v1}}
\item {\scriptsize\url{https://openai.com/index/dall-e-api-now-available-in-public-beta}}
\item {\scriptsize\url{https://openai.com/index/dall-e-now-available-in-beta}}
\item {\scriptsize\url{https://openai.com/index/dall-e-now-available-without-waitlist}}
\item {\scriptsize\url{https://openai.com/index/hierarchical-text-conditional-image-generation-with-clip-latents}}
\end{itemize}
\endgroup
\end{samepage}
\end{minipage}
\endgroup
\end{technicalnote}
\medskip

\begin{technicalnote}{Stable Diffusion}{主要資料}
\begingroup
% reader-facing Technical Notes break policy
\widowpenalty=10000
\clubpenalty=10000
\displaywidowpenalty=10000
\begin{tabularx}{\linewidth}{@{}>{\bfseries}p{0.22\linewidth}X@{}}
組織 & Stability AI / collaborators \\
種別 & オープンウェイト \\
時系列 & — \\
\end{tabularx}
\smallskip
{\bfseries 一次資料から整理したtechnical points}
\begin{itemize}[leftmargin=1.5em,itemsep=0.35em]
\item \textbf{一次情報で確認できる事実}: Stability AI / collaboratorsの選定済み一次資料では、Stability AIは2022年8月10日にStable Diffusionをまずresearcher向けにreleaseし、model weightsはaccess取得後にHugging Faceで提供、codeとmodel cardも公開する段階的releaseとして位置付けた。 8月22日のpublic releaseではweights・model card・codeを一般公開し、Creative ML OpenRAIL-M licenseの下でcommercial/non-commercial利用を認める公開境界へ移行した。
\end{itemize}
\begin{minipage}{\linewidth}
% reader-facing Technical Notes generic boundary/source tail group
\begin{samepage}
{\bfseries 一次資料}
\begingroup\sloppy
\Urlmuskip=0mu plus 2mu\relax
\begin{itemize}[leftmargin=1.5em,itemsep=0.25em]
\item {\scriptsize\url{https://stability.ai/news-updates/stable-diffusion-announcement}}
\item {\scriptsize\url{https://stability.ai/news-updates/stable-diffusion-public-release}}
\end{itemize}
\endgroup
\end{samepage}
\end{minipage}
\endgroup
\end{technicalnote}
\medskip

\begin{technicalnote}{Imagen}{補足資料}
\begingroup
% reader-facing Technical Notes break policy
\widowpenalty=10000
\clubpenalty=10000
\displaywidowpenalty=10000
\begin{tabularx}{\linewidth}{@{}>{\bfseries}p{0.22\linewidth}X@{}}
組織 & Google Research / authors \\
種別 & 論文 \\
時系列 & — \\
\end{tabularx}
\smallskip
{\bfseries 一次資料から整理したtechnical points}
\begin{itemize}[leftmargin=1.5em,itemsep=0.35em]
\item \textbf{一次情報で確認できる事実}: Google Research / authorsの選定済み一次資料では、対象event近傍の一次資料から image generation を確認できる。
\end{itemize}
\begin{minipage}{\linewidth}
% reader-facing Technical Notes generic boundary/source tail group
\begin{samepage}
{\bfseries 一次資料}
\begingroup\sloppy
\Urlmuskip=0mu plus 2mu\relax
\begin{itemize}[leftmargin=1.5em,itemsep=0.25em]
\item {\scriptsize\url{http://arxiv.org/abs/2205.11487v1}}
\end{itemize}
\endgroup
\end{samepage}
\end{minipage}
\endgroup
\end{technicalnote}
\medskip

\begin{technicalnote}{LAION-5B}{補足資料}
\begingroup
% reader-facing Technical Notes break policy
\widowpenalty=10000
\clubpenalty=10000
\displaywidowpenalty=10000
\begin{tabularx}{\linewidth}{@{}>{\bfseries}p{0.22\linewidth}X@{}}
組織 & LAION / authors \\
種別 & 論文 \\
時系列 & — \\
\end{tabularx}
\smallskip
{\bfseries 一次資料から整理したtechnical points}
\begin{itemize}[leftmargin=1.5em,itemsep=0.35em]
\item \textbf{一次情報で確認できる事実}: LAION / authorsの選定済み一次資料では、対象event近傍の一次資料から 5B parameter scale を確認できる。
\end{itemize}
\begin{minipage}{\linewidth}
% reader-facing Technical Notes generic boundary/source tail group
\begin{samepage}
{\bfseries 一次資料}
\begingroup\sloppy
\Urlmuskip=0mu plus 2mu\relax
\begin{itemize}[leftmargin=1.5em,itemsep=0.25em]
\item {\scriptsize\url{http://arxiv.org/abs/2210.08402v1}}
\end{itemize}
\endgroup
\end{samepage}
\end{minipage}
\endgroup
\end{technicalnote}
\medskip

\begin{technicalnote}{DreamBooth}{補足資料}
\begingroup
% reader-facing Technical Notes break policy
\widowpenalty=10000
\clubpenalty=10000
\displaywidowpenalty=10000
\begin{tabularx}{\linewidth}{@{}>{\bfseries}p{0.22\linewidth}X@{}}
組織 & Google Research / authors \\
種別 & 論文 \\
時系列 & — \\
\end{tabularx}
\smallskip
{\bfseries 一次資料から整理したtechnical points}
\begin{itemize}[leftmargin=1.5em,itemsep=0.35em]
\item \textbf{一次情報で確認できる事実}: Google Research / authorsの選定済み一次資料では、DreamBoothは少数のsubject画像を入力としてpretrained text-to-image diffusion modelをfine-tuneし、unique identifierを特定subjectへbindingするpersonalization手法として提案された。 著者らはclass-specific prior preservation lossを導入し、reference画像そのものにないscene・pose・view・lightingへsubjectを再文脈化する設計を扱った。
\end{itemize}
\begin{minipage}{\linewidth}
% reader-facing Technical Notes generic boundary/source tail group
\begin{samepage}
{\bfseries 一次資料}
\begingroup\sloppy
\Urlmuskip=0mu plus 2mu\relax
\begin{itemize}[leftmargin=1.5em,itemsep=0.25em]
\item {\scriptsize\url{http://arxiv.org/abs/2208.12242v2}}
\end{itemize}
\endgroup
\end{samepage}
\end{minipage}
\endgroup
\end{technicalnote}
\medskip

\begin{technicalnote}{AudioLM}{補足資料}
\begingroup
% reader-facing Technical Notes break policy
\widowpenalty=10000
\clubpenalty=10000
\displaywidowpenalty=10000
\begin{tabularx}{\linewidth}{@{}>{\bfseries}p{0.22\linewidth}X@{}}
組織 & Google Research / authors \\
種別 & 論文 \\
時系列 & — \\
\end{tabularx}
\smallskip
{\bfseries 一次資料から整理したtechnical points}
\begin{itemize}[leftmargin=1.5em,itemsep=0.35em]
\item \textbf{一次情報で確認できる事実}: Google Research / authorsの選定済み一次資料では、AudioLMはaudioをdiscrete token列へ写像し、そのrepresentation space上のlanguage modelingとしてaudio generationを定式化した。 長期構造を捉えるmasked-language-model由来tokenと、再構成品質を担うneural audio codec由来codeを組み合わせるhybrid tokenizationを採用した。
\end{itemize}
\begin{minipage}{\linewidth}
% reader-facing Technical Notes generic boundary/source tail group
\begin{samepage}
{\bfseries 一次資料}
\begingroup\sloppy
\Urlmuskip=0mu plus 2mu\relax
\begin{itemize}[leftmargin=1.5em,itemsep=0.25em]
\item {\scriptsize\url{http://arxiv.org/abs/2209.03143v2}}
\end{itemize}
\endgroup
\end{samepage}
\end{minipage}
\endgroup
\end{technicalnote}
\medskip

\begin{technicalnote}{Make-A-Video}{補足資料}
\begingroup
% reader-facing Technical Notes break policy
\widowpenalty=10000
\clubpenalty=10000
\displaywidowpenalty=10000
\begin{tabularx}{\linewidth}{@{}>{\bfseries}p{0.22\linewidth}X@{}}
組織 & Meta AI / authors \\
種別 & 論文 \\
時系列 & — \\
\end{tabularx}
\smallskip
{\bfseries 一次資料から整理したtechnical points}
\begin{itemize}[leftmargin=1.5em,itemsep=0.35em]
\item \textbf{一次情報で確認できる事実}: Meta AI / authorsの選定済み一次資料では、Make-A-Videoはtext-image pairから外観と言語対応を学び、unlabeled videoからmotionを学ぶことで、paired text-video dataを必須としないtext-to-video生成を構成した。 text-to-image modelへspatio-temporal moduleを追加し、video decoder・interpolation・super-resolutionを組み合わせる生成pipelineを設計した。
\end{itemize}
\begin{minipage}{\linewidth}
% reader-facing Technical Notes generic boundary/source tail group
\begin{samepage}
{\bfseries 一次資料}
\begingroup\sloppy
\Urlmuskip=0mu plus 2mu\relax
\begin{itemize}[leftmargin=1.5em,itemsep=0.25em]
\item {\scriptsize\url{http://arxiv.org/abs/2209.14792v1}}
\end{itemize}
\endgroup
\end{samepage}
\end{minipage}
\endgroup
\end{technicalnote}
\medskip

\begin{technicalnote}{DreamFusion}{補足資料}
\begingroup
% reader-facing Technical Notes break policy
\widowpenalty=10000
\clubpenalty=10000
\displaywidowpenalty=10000
\begin{tabularx}{\linewidth}{@{}>{\bfseries}p{0.22\linewidth}X@{}}
組織 & Google Research / authors \\
種別 & 論文 \\
時系列 & — \\
\end{tabularx}
\smallskip
{\bfseries 一次資料から整理したtechnical points}
\begin{itemize}[leftmargin=1.5em,itemsep=0.35em]
\item \textbf{一次情報で確認できる事実}: Google Research / authorsの選定済み一次資料では、DreamFusionはpretrained 2D text-to-image diffusion modelをpriorとして利用し、large-scale 3D training dataを直接必要とせずtext-to-3D synthesisを行う方法を提案した。 probability density distillationに基づくlossを用い、randomly initialized NeRFの複数視点renderingをgradient descentで最適化する構成を採る。
\end{itemize}
\begin{minipage}{\linewidth}
% reader-facing Technical Notes generic boundary/source tail group
\begin{samepage}
{\bfseries 一次資料}
\begingroup\sloppy
\Urlmuskip=0mu plus 2mu\relax
\begin{itemize}[leftmargin=1.5em,itemsep=0.25em]
\item {\scriptsize\url{http://arxiv.org/abs/2209.14988v1}}
\end{itemize}
\endgroup
\end{samepage}
\end{minipage}
\endgroup
\end{technicalnote}
\medskip
